% ======================================================================
\section{Conclusion}\label{sec:conclusion}
% ======================================================================

The standard analysis of SGD rests on a one-step expected decrease of the function value.
For SGD with momentum the function value is the wrong yardstick, because momentum makes
the iterates overshoot. The report's main message fits in three sentences.

\begin{enumerate}
\item \emph{SGDM is SGD on the extrapolated points.} The points
  $z_k=x_k+\frac{\beta}{1-\beta}(x_k-x_{k-1})$, where the ball would come to rest if the
  force were switched off, follow exact SGD steps $z_{k+1}=z_k-\gamma g_k$ with the effective
  step size $\gamma=\alpha/(1-\beta)$; only the gradient is evaluated at the nearby point $x_k$
  (\cref{lem:z}).
\item \emph{A potential decreases in expectation, exactly as $f$ does for SGD.} The
  potential $\Phi_k=f(z_k)-\fs+\frac{L\beta^2}{2(1-\beta)^2}\norm{x_k-x_{k-1}}^2$ satisfies
  $\Ek\Phi_{k+1}\le\Phi_k-\frac\gamma4\norm{\nabla f(x_k)}^2+L\gamma^2\sigma^2$ whenever
  $\alpha\le(1-\beta)^2/(2L)$: friction must dissipate kinetic energy faster than the gradient
  mismatch creates error (\cref{thm:master}).
\item \emph{Everything else follows as for SGD.} Telescoping gives the $O(1/\sqrt K)$
  nonconvex rate, diminishing steps give asymptotic convergence, the PL inequality gives
  linear convergence to a noise floor, and convexity relaxes the step-size condition to
  $\gamma\le1/(2L)$ (\cref{sec:convergence}).
\end{enumerate}

As for \emph{when} convergence holds (\cref{sec:circumstances}): the momentum must satisfy
$\beta<1$; the step size must be small in terms of the effective step $\gamma$, with a
threshold that depends on the function class (\cref{fig:regions}); a constant step only
reaches a noise floor of order $\gamma\sigma^2$, which is unavoidable and the same as for SGD
(\cref{prop:noise-floor}); and stability or optimality on quadratics does not imply
convergence on general functions (\cref{ex:lrp}). Once step sizes are compared on the
right scale, our upper bounds for SGDM match those for SGD in the dominant noise term, up to
constants, and they show no acceleration.

The potential is a convenience, not a necessity: \cref{sec:lyapunov-free} reaches the same
nonconvex rate by tracking only $f(x_k)$, at the price of the slightly stronger restriction
$\beta\le1-4L\gamma$ on the momentum. There, momentum is an exponentially averaged gradient
estimate whose bias is controlled by the slow motion of the iterates and whose averaged noise
is almost orthogonal to the current gradient. \Cref{sec:lf-all} extends this argument to
the other settings of \cref{sec:convergence}: exact gradients, decreasing steps and the PL
inequality with decaying steps (both in the momentum-buffer form used in practice, which the
potential argument does not cover for varying steps), the PL inequality with a constant
step, and convex functions (where, as for SGD, the distance to a minimizer is tracked); it
also handles noise that grows with the gradient. One weighted error budget
(\cref{lem:weighted-budget}) does most of the work. What the potential buys is flexibility.
The Lyapunov-free argument needs the iterates to move slowly compared with the memory of the
momentum, $L\gamma\lesssim1-\beta$, whereas the energy argument of \cref{prop:energy} and the
convex \cref{thm:convex} allow effective steps that do not shrink as $\beta\to1$; and with
varying steps it needs a non-increasing schedule fixed in advance.
